\documentclass[11pt]{article}
\usepackage[a4paper,margin=25mm]{geometry}
\usepackage{amsmath,amssymb,amsthm,hyperref}
\newtheorem{theorem}{Theorem}
\newcommand{\D}[1]{\lvert D_{16}^{(#1)}\rangle}
\newcommand{\R}{\mathcal R}
\title{An analytical third-order anticoherent plane for sixteen qubits}
\author{Catalogue calculation}
\date{30 October 2026}
\begin{document}
\maketitle
Throughout, $\D r$ is the normalized symmetric Dicke state with $r$ excitations.
The reduced operator $\R_t$ traces out $16-t$ qubits and is represented in
normalized Dicke coordinates on the remaining symmetric space.
A subspace is $t$-anticoherent when every normalized vector in it has
$\R_t(\lvert\psi\rangle\langle\psi\rvert)=I_{t+1}/(t+1)$.

\section{Explicit normalized basis}
Define three unit complex numbers
\begin{equation}\label{eq:phases}
z=\frac{545-192i\sqrt{10}}{7\sqrt{13585}},\qquad
f=-\frac{55+4i\sqrt{10}}{7\sqrt{65}},\qquad
g=\frac{19+8i\sqrt{10}}{\sqrt{1001}}.
\end{equation}
Their moduli are one, as follows immediately by squaring their real and
imaginary parts. Set
\begin{align}
u_1&=\frac{\sqrt{38}}{16}\bigl(\D0+z^2\D{16}\bigr)
 +\frac z8\bigl(\sqrt{10}\D4+5\D8+\sqrt{10}\D{12}\bigr),\label{eq:u1}\\
u_2&=\frac{\sqrt5}{4}\bigl(\D2+fg\D{14}\bigr)
 +\frac{\sqrt3}{4}\bigl(f\D6+g\D{10}\bigr).\label{eq:u2}
\end{align}
\begin{theorem}
The vectors $u_1,u_2$ are orthonormal and satisfy
\begin{equation}\label{eq:cross}
\R_3(\lvert u_a\rangle\langle u_b\rvert)=\delta_{ab}\frac{I_4}{4}
\qquad(a,b\in\{1,2\}).
\end{equation}
Consequently their complex linear span is a two-dimensional third-order
anticoherent subspace. Its dimension is maximal: $K(16,3)=2$.
\end{theorem}

\section{Full-span verification}
The squared Dicke amplitudes, in increasing order on the stated supports, are
\begin{equation}
 u_1:\quad \frac1{128}(19,20,50,20,19),\qquad
 u_2:\quad \frac1{16}(5,3,3,5).
\end{equation}
Each list sums to one, and the supports are disjoint, proving
orthonormality. To check all reduced operators explicitly, associate to
$v=\sum_{r=0}^{16}v_r\D r$ its bipartition matrix
\begin{equation}\label{eq:matrix}
 [M_t(v)]_{a,b}=v_{a+b}
 \sqrt{\frac{\binom ta\binom{16-t}b}{\binom{16}{a+b}}},
 \quad 0\le a\le t,\quad0\le b\le16-t.
\end{equation}
Then $\R_t(\lvert v\rangle\langle w\rvert)=M_t(v)M_t(w)^\dagger$.
Within each vector, consecutive occupied Dicke indices differ by four.
For $t=3$ all off-diagonal entries of both self reductions therefore
vanish. Substitution of the two probability lists into
\eqref{eq:matrix} gives
\begin{equation}
 M_3(u_1)M_3(u_1)^\dagger=M_3(u_2)M_3(u_2)^\dagger
 =\operatorname{diag}(1/4,1/4,1/4,1/4).
\end{equation}

It remains essential to check the cross reductions: separate
anticoherence of the displayed basis vectors alone would not suffice.
Write $F=\overline f$, $G=\overline g$, and put
\begin{align}
 E_1&=7\sqrt{19}+z\bigl(\sqrt{91}FG+15\sqrt{11}F+9\sqrt{35}G\bigr),\\
 E_2&=\sqrt{91}FG+\sqrt{11}F+2\sqrt{35}G,\\
 E_3&=2\sqrt{35}F+\sqrt{11}G+\sqrt{91},\\
 E_4&=7\sqrt{19}FGz+9\sqrt{35}F+15\sqrt{11}G+\sqrt{91}.
\end{align}
Direct substitution of \eqref{eq:phases} gives
$E_1=E_2=E_3=E_4=0$, with exact radical arithmetic. All possible nonzero
entries of the first cross reduction are displayed in
\begin{equation}
 M_3(u_1)M_3(u_2)^\dagger=
 \begin{pmatrix}
 0&0&E_1/896&0\\
 0&0&0&3zE_2/448\\
 3zE_3/448&0&0&0\\
 0&zE_4/896&0&0
 \end{pmatrix}=0.
\end{equation}
The reverse ordered cross reduction is its adjoint. This proves
\eqref{eq:cross}, including every off-diagonal basis pair. For arbitrary
complex $\alpha,\beta$ with $|\alpha|^2+|\beta|^2=1$, linearity now yields
\begin{equation}
 \R_3\bigl(\lvert\alpha u_1+\beta u_2\rangle
 \langle\alpha u_1+\beta u_2\rvert\bigr)=I_4/4.
\end{equation}

The two basis vectors individually have exact AC order three. Indeed,
their fourth-reduction purities are
\begin{equation}
 \operatorname{Tr}\R_4(\lvert u_1\rangle\langle u_1\rvert)^2
 =\frac{3457}{16562}>\frac15,\qquad
 \operatorname{Tr}\R_4(\lvert u_2\rangle\langle u_2\rvert)^2
 =\frac{3433}{16562}>\frac15.
\end{equation}
In particular, the whole plane is not fourth-order anticoherent.
These two purity computations do not assert that every superposition
has exactly order three.

\section{Why a larger third-order subspace is impossible}
Here is the relevant general dimension bound, specialized so that the
optimality statement is self-contained. Let
$\mathcal H_N=\operatorname{Sym}^N(\mathbb C^2)$ and let
$V\subset\mathcal H_{16}$ be an arbitrary third-order AC subspace,
with projector $P$ and dimension $k$.
The polarized AC identities are
\begin{equation}
 PTP=\frac{\operatorname{Tr}T}{17}P
\end{equation}
for every operator $T$ having spin-tensor ranks at most three.
Choose Hilbert--Schmidt orthonormal rectangular irreducible tensor maps
$A_{\ell m}:\mathcal H_{16}\to\mathcal H_{15}$ with
$\ell=1/2,3/2$. These exist by the spin decomposition of
$\operatorname{Hom}(\mathcal H_{16},\mathcal H_{15})$.
There are $2+4=6$ such maps. Each product
$A_{\ell m}^{\dagger}A_{\ell'm'}$ has tensor ranks at most
$\ell+\ell'\le3$, and its scalar part is
$\delta_{\ell\ell'}\delta_{mm'}I_{17}/17$.
For an orthonormal basis $v_1,\ldots,v_k$ of $V$, it follows that
\begin{equation}
 \langle A_{\ell m}v_a,A_{\ell'm'}v_b\rangle
 =\frac{\delta_{ab}\delta_{\ell\ell'}\delta_{mm'}}{17}.
\end{equation}
Thus $6k$ nonzero, mutually orthogonal vectors lie in the
16-dimensional space $\mathcal H_{15}$, so $6k\le16$ and $k\le2$.
Together with \eqref{eq:u1}--\eqref{eq:u2}, this proves $K(16,3)=2$.
The argument applies to arbitrary complex subspaces, with no support
or symmetry assumption.

\section{Reconstruction and scope}
The sparse supports and the simple probability lists above were obtained
from the numerical example \texttt{\(\Pi\)\_16\_2\_3} in the supplied
\texttt{ACLiege.jl}, at line 901. Fixing these probabilities reduces the
remaining full-span conditions to the four phase identities displayed
above; their solution \eqref{eq:phases} gives an exact certificate.
A common collective rotation $\exp(-i\beta J_z)$, with
$\beta\simeq-0.7306438994490132$, and independent phases on the two
basis vectors match the supplied numerical vectors to maximum Dicke
amplitude errors $3.32\times10^{-10}$ and $5.10\times10^{-10}$.
This comparison is approximate and is separate from the exact proof.

The machine-readable \texttt{basis.json} contains every exact and
numerical Dicke coefficient. The primary and independent verifiers
rebuild all 64 scalar entries of the four ordered cross reductions in
exact rational and radical arithmetic. The saved reports also record
the fourth-reduction purities and the source comparison. No claim of
classification of all optimal subspaces or of literature priority is made.
\end{document}
